\documentclass[../../main.tex]{subfiles}
\begin{document}
\section{流形的定向}\label{section:ch15-manifolds}

设 $M$ 是一个带边或不带边的光滑流形。我们把 $M$ 上的一个\emph{逐点定向}定义为对每个切空间定向的一个选取。就其本身而言，这并不是一个很有用的概念，因为邻近各点的定向之间可能毫无关系。例如，$\mathbb{R}^n$ 上的一个逐点定向可能在每一点处随机地在标准定向与其相反定向之间切换。为使定向与光滑结构产生某种联系，我们需要一个额外条件，以保证邻近切空间的定向彼此一致。

设 $M$ 是一个带边或不带边的光滑 $n$ 流形，赋予了一个逐点定向。若 $(E_i)$ 是 $TM$ 的一个局部标架，我们说 $(E_i)$ 是（正）定向的，如果在每一点 $p\in U$ 处 $(E_1|_p,\dots,E_n|_p)$ 都是 $T_pM$ 的一组正定向基。负定向标架类似地定义。

若 $M$ 的每一点都含于某个定向局部标架的定义域内，就称该逐点定向是\emph{连续的}。（回忆一下，由定义，构成一个局部标架的那些向量场是连续的。）$M$ 的一个定向就是一个连续的逐点定向。若 $M$ 存在一个定向，就说它是\emph{可定向的}，否则说它不可定向。一个\emph{有向流形}是一个有序对 $(M,\mathcal O)$，其中 $M$ 是一个可定向光滑流形，$\mathcal O$ 是对 $M$ 的一个定向的选取；带边有向流形类似地定义。对每个 $p\in M$，由 $\mathcal O$ 确定的 $T_pM$ 的定向记作 $\mathcal O_p$。当不必显式命名该定向时，我们使用通常的简写表达“$M$ 是一个有向光滑流形”（或“带边流形”）。

若 $M$ 是零维的，这个定义只不过意味着 $M$ 的一个定向是附加在它每一点上的一个 $\pm1$ 的选取。在此情形下连续性条件是空洞的，定向标架的概念也无用处。显然，每个 $0$-流形都是可定向的。

\begin{exercise}\label{exercise:ch15-04}
设 $M$ 是一个带边或不带边的有向光滑 $n$ 流形，$n\geqslant1$。证明：定义域连通的每个局部标架要么是正定向的，要么是负定向的。证明连通性假设是必要的。
\end{exercise}

下面两个命题给出比定义更便于使用的指定流形定向的方法。

\begin{proposition}[$n$-形式所确定的定向]\label{proposition:ch15-nform-orientation}
设 $M$ 是一个带边或不带边的光滑 $n$ 流形。$M$ 上任意一个处处非零的 $n$-形式 $\omega$ 都确定 $M$ 的唯一一个定向，使得 $\omega$ 在每一点都正定向。反之，若 $M$ 被赋予一个定向，则 $M$ 上存在一个光滑、处处非零、且在每一点都正定向的 $n$-形式。
\end{proposition}

\begin{remark}
由于这个命题，若 $M$ 是一个带边或不带边的光滑 $n$ 流形，$M$ 上任意处处非零的 $n$-形式都称为一个\emph{定向形式}。若 $M$ 已定向，$\omega$ 是确定该给定定向的一个定向形式，我们也说 $\omega$ 是（正）定向的。容易验证，若 $\omega$ 与 $\widetilde\omega$ 是 $M$ 上两个正定向的光滑形式，则 $\widetilde\omega=f\omega$，其中 $f$ 是某个严格正的光滑实值函数。若 $M$ 是一个 $0$-流形，一个处处非零的 $0$-形式（即一个实值函数）在其为正的点处指派定向 $+1$，在其为负的点处指派定向 $-1$。
\end{remark}

\begin{proof}
设 $\omega$ 是 $M$ 上一个处处非零的 $n$-形式。则 $\omega$ 由命题\ref{proposition:ch15-n-covector-orientation}定义了一个逐点定向，故我们只需验证它连续。当 $n=0$ 时这平凡成立，故设 $n\geqslant1$。给定 $p\in M$，在 $p$ 的一个连通邻域 $U$ 上取任意局部标架 $(E_i)$，令 $(\varepsilon^i)$ 为对偶余标架。在 $U$ 上，$\omega$ 在此标架中的表达式为 $\omega=f\,\varepsilon^1\wedge\cdots\wedge\varepsilon^n$，其中 $f$ 是某个连续函数。$\omega$ 处处非零意味着 $f$ 处处非零，因而
\begin{align*}
\omega(E_1,\dots,E_n)=f\ne0
\end{align*}
在 $U$ 的所有点处成立。由于 $U$ 连通，这个表达式在 $U$ 上要么恒正，要么恒负，因而给定标架要么正定向，要么负定向。若为负定向，我们可以把 $E_1$ 换成 $-E_1$，从而得到一个正定向的新标架。于是由 $\omega$ 确定的逐点定向是连续的。

反之，设 $M$ 已定向，令 $\Lambda^n_+T^*M\subseteq\Lambda^n T^*M$ 为 $M$ 各点处由正定向 $n$-余向量组成的开子集。在任意点 $p\in M$ 处，$\Lambda^n_+T^*M$ 与纤维 $\Lambda^n(T^*_pM)$ 的交是一条开半直线，因而是凸的。由通常的单位分解论证（见习题\ref{exercise:ch13-02}），存在 $\Lambda^n_+T^*M$ 的一个光滑整体截面（即一个正定向的光滑整体 $n$-形式）。
\end{proof}

在一个带边或不带边的有向光滑流形上，若坐标标架 $(\partial/\partial x^i)$ 正定向，就称该光滑坐标图是（正）定向的；若坐标标架负定向，就称它是负定向的。若对每对 $\alpha,\beta$，过渡映射 $\varphi_\beta\circ\varphi_\alpha^{-1}$ 在 $\varphi_\alpha(U_\alpha\cap U_\beta)$ 上处处都有正的 Jacobian 行列式，就称光滑坐标图册 $\{(U_\alpha,\varphi_\alpha)\}$ 是\emph{一致定向的}。

\begin{proposition}[坐标图册所确定的定向]\label{proposition:ch15-atlas-orientation}
设 $M$ 是一个带边或不带边的正维光滑流形。给定 $M$ 的任意一个一致定向的光滑图册，存在 $M$ 的唯一一个定向，使给定图册中每个坐标图都正定向。反之，若 $M$ 已定向且或者 $\partial M=\varnothing$，或者 $\dim M>1$，则全体定向光滑坐标图所成之集是 $M$ 的一个一致定向图册。
\end{proposition}

\begin{proof}
首先，设 $M$ 有一个一致定向的光滑图册。图册中每个坐标图都在其定义域的每一点处确定一个逐点定向。当两个坐标图重叠时，它们各自坐标标架之间的过渡矩阵就是过渡映射的 Jacobian 矩阵，由假设它有正行列式，故它们在每一点处确定相同的逐点定向。这样确定的定向是连续的，因为每一点都含于某个定向坐标标架的定义域内。

反之，设 $M$ 已定向且或者 $\partial M=\varnothing$，或者 $\dim M>1$。每一点都含于某个光滑坐标图的定义域内；若该坐标图负定向，我们可以把 $x^1$ 换成 $-x^1$，从而得到一个正定向的新坐标图。这些坐标图都正定向这一事实保证了它们的过渡映射有正的 Jacobian 行列式，故它们构成一个一致定向图册。（当 $\dim M=1$ 时，由于我们在边界坐标图中约定最后一个坐标非负，这一做法对边界坐标图行不通。）
\end{proof}

\begin{proposition}[乘积定向]\label{proposition:ch15-product-orientation}
设 $M_1,\dots,M_k$ 是可定向光滑流形。在 $M_1\times\cdots\times M_k$ 上存在唯一一个定向，称为\emph{乘积定向}，具有下列性质：若对每个 $i=1,\dots,k$，$\omega_i$ 是 $M_i$ 上给定定向的一个定向形式，则 $\pi_1^*\omega_1\wedge\cdots\wedge\pi_k^*\omega_k$ 是乘积定向的一个定向形式。
\end{proposition}

\begin{exercise}\label{exercise:ch15-08}
证明上一命题。
\end{exercise}

\begin{proposition}\label{proposition:ch15-two-orientations}
设 $M$ 是一个连通的、可定向的、带边或不带边的光滑流形。则 $M$ 恰有两个定向。若 $M$ 的两个定向在某一点处一致，则它们相等。
\end{proposition}

\begin{exercise}\label{exercise:ch15-10}
证明上一命题。
\end{exercise}

\begin{proposition}[余维 $0$ 子流形的定向]\label{proposition:ch15-codim0-orientation}
设 $M$ 是一个带边或不带边的有向光滑流形，$D\subseteq M$ 是一个带边或不带边的光滑余维 $0$ 子流形。则 $M$ 的定向限制为 $D$ 的一个定向。若 $\omega$ 是 $M$ 的一个定向形式，则 $\iota_D^*\omega$ 是 $D$ 的一个定向形式。
\end{proposition}

\begin{exercise}\label{exercise:ch15-12}
证明上一命题。
\end{exercise}

设 $M,N$ 是带边或不带边的有向光滑流形，$F:M\to N$ 是一个局部微分同胚。若 $M,N$ 都是正维的，当对每个 $p\in M$，同构 $\mathrm{d}F_p$ 把 $T_pM$ 的定向基映到 $T_{F(p)}N$ 的定向基时，就说 $F$ 是\emph{保定向的}；当它把 $T_pM$ 的定向基映到 $T_{F(p)}N$ 的负定向基时，就说 $F$ 是\emph{反定向的}。若 $M,N$ 都是 $0$-流形，则当对每个 $p\in M$，点 $p$ 与 $F(p)$ 有相同时定向时 $F$ 保定向；当它们有相反定向时 $F$ 反定向。

\begin{exercise}\label{exercise:ch15-13}
设 $M,N$ 是带边或不带边的正维有向光滑流形，$F:M\to N$ 是一个局部微分同胚。证明下列各条等价。
\begin{enumerate}[(a)]
\item $F$ 保定向。
\item 关于 $M,N$ 的任意定向光滑坐标图，$F$ 的 Jacobian 矩阵行列式为正。
\item 对 $N$ 的任意正定向定向形式 $\omega$，形式 $F^*\omega$ 在 $M$ 上正定向。
\end{enumerate}
\end{exercise}

\begin{exercise}\label{exercise:ch15-14}
证明保定向映射的复合仍保定向。
\end{exercise}

下面是另一个构造定向的重要方法。

\begin{proposition}[拉回定向]\label{proposition:ch15-pullback-orientation}
设 $M,N$ 是带边或不带边的光滑流形。若 $F:M\to N$ 是一个局部微分同胚且 $N$ 可定向，则 $M$ 有唯一一个定向，称为由 $F$ 诱导的\emph{拉回定向}，使 $F$ 保定向。
\end{proposition}

\begin{proof}
对每个 $p\in M$，在 $T_pM$ 上存在唯一一个定向，使同构 $\mathrm{d}F_p:T_pM\to T_{F(p)}N$ 保定向。这就在 $M$ 上定义了一个逐点定向；只要它连续，它就是使 $F$ 保定向的 $M$ 上唯一的定向。为看清它连续，只需为 $N$ 选取一个光滑定向形式 $\omega$，并注意 $F^*\omega$ 是 $M$ 上的一个光滑定向形式。
\end{proof}

在上一命题的情形中，若 $\mathcal O$ 表示 $N$ 上给定的定向，则 $M$ 上的拉回定向记作 $F^*\mathcal O$。

\begin{exercise}\label{exercise:ch15-16}
设 $F:M\to N$、$G:N\to P$ 是局部微分同胚，$\mathcal O$ 是 $P$ 上的一个定向。证明 $(G\circ F)^*\mathcal O=F^*(G^*\mathcal O)$。
\end{exercise}

回忆一下，若一个光滑流形容许一个光滑整体标架，就称它是可平行化的。

\begin{proposition}\label{proposition:ch15-parallelizable-orientable}
每个可平行化的光滑流形都是可定向的。
\end{proposition}

\begin{proof}
设 $M$ 可平行化，令 $(E_1,\dots,E_n)$ 是 $M$ 的一个整体光滑标架。通过宣布基 $(E_1|_p,\dots,E_n|_p)$ 在每个 $p\in M$ 处正定向，来在 $M$ 上定义一个逐点定向。这个逐点定向是连续的，因为 $M$ 的每一点都含于（整体）定向标架 $(E_i)$ 的定义域内。
\end{proof}

\begin{figure}[htbp]
\centering
\includegraphics[width=0.7\textwidth]{fig15-4-vectorfield.png}
\caption{沿子流形的一个向量场}
\label{figure:ch15-vectorfield}
\end{figure}

\begin{example}\label{example:ch15-18}
上一命题表明 Euclid 空间 $\mathbb{R}^n$、$n$-环面 $\mathbb T^n$、球面 $S^1,S^3,S^7$ 以及它们的乘积都是可定向的，因为它们都可平行化。因此，这些流形中任意一个的余维 $0$ 子流形也可定向。类似地，每个 Lie 群都可定向，因为它可平行化。
\end{example}

在 Lie 群的情形，我们还能说得更多。若 $G$ 是一个 Lie 群，当 $L_g$ 对每个 $g\in G$ 都保定向时，就说 $G$ 的一个定向是\emph{左不变的}。

\begin{proposition}\label{proposition:ch15-lie-group-orientation}
每个 Lie 群恰有两个左不变定向，分别对应于其 Lie 代数的两个定向。
\end{proposition}

\begin{exercise}\label{exercise:ch15-20}
证明上一命题。
\end{exercise}

\subsection{超曲面的定向}

若 $M$ 是一个有向光滑流形，$S$ 是 $M$ 的一个光滑子流形（带边或不带边），$S$ 未必从 $M$ 继承一个定向；即使 $S$ 嵌入也是如此。显然，把一个定向形式从 $M$ 限制到 $S$ 是不够的，因为一个 $n$-形式限制到一个低维流形上必然为零。一个有用的例子是 Möbius 带，它不可定向（见例\ref{example:ch15-mobius-bundle}），尽管它可嵌入 $\mathbb{R}^3$。

本节中我们把注意力集中于浸入或嵌入的超曲面（余维 $1$ 子流形）。借助额外的一条信息（一个处处不与超曲面相切的向量场），我们可以用 $M$ 上的一个定向来诱导 $M$ 中超曲面上的一个定向。

设 $M$ 是一个带边或不带边的光滑流形，$S\subseteq M$ 是一个光滑子流形（浸入或嵌入，带边或不带边）。回忆一下（例\ref{example:ch10-10}），沿 $S$ 的一个向量场是 ambient 切丛 $TM|_S$ 的一个截面，即一个连续映射 $N:S\to TM$，满足对每个 $p\in S$ 有 $N_p\in T_pM$（图\ref{figure:ch15-vectorfield}）。例如，$M$ 上任意向量场都限制为沿 $S$ 的一个向量场；但一般而言，并非沿 $S$ 的每个向量场均如此（见习题\ref{exercise:ch10-09}）。

\begin{figure}[htbp]
\centering
\includegraphics[width=0.7\textwidth]{fig15-5-induced.png}
\caption{由一个处处不相切的向量场所诱导的定向}
\label{figure:ch15-induced}
\end{figure}

\begin{proposition}\label{proposition:ch15-hypersurface-orientation}
设 $M$ 是一个带边或不带边的有向光滑 $n$ 流形，$S$ 是 $M$ 中一个带边或不带边的浸入超曲面，$N$ 是沿 $S$ 的一个处处不与 $S$ 相切的向量场。则 $S$ 有唯一一个定向，使对每个 $p\in S$，$(E_1,\dots,E_{n-1})$ 是 $T_pS$ 的一个定向基当且仅当 $(N_p,E_1,\dots,E_{n-1})$ 是 $T_pM$ 的一个定向基。若 $\omega$ 是 $M$ 的一个定向形式，则关于这个定向，$\iota_S^*(N\lrcorner\omega)$ 是 $S$ 的一个定向形式，其中 $\iota_S:S\hookrightarrow M$ 是包含映射。
\end{proposition}

\begin{remark}
$n=3$ 的情形见图\ref{figure:ch15-induced}。当 $n=1$ 时，由于 $S$ 是一个 $0$-流形，这个命题应解释如下：在每个 $p\in S$ 处，若 $N_p$ 是 $T_pM$ 的一个定向基，则我们给 $p$ 指派定向 $+1$；若 $N_p$ 负定向，则指派 $-1$。在这样理解下，下面的证明在 $n=1$ 的情形无需修改即可通过。
\end{remark}

\begin{proof}
设 $\omega$ 是 $M$ 的一个定向形式。则 $\Theta=\iota_S^*(N\lrcorner\omega)$ 是 $S$ 上的一个 $(n-1)$-形式。（回忆一下，拉回 $\iota_S^*$ 实际上就是限制到与 $S$ 相切的向量上。）若我们能证明 $\Theta$ 处处不为零，便知它是 $S$ 的一个定向形式。给定 $T_pS$ 的任意基 $(E_1,\dots,E_{n-1})$，$N$ 处处不与 $S$ 相切这一事实意味着 $(N_p,E_1,\dots,E_{n-1})$ 是 $T_pM$ 的一个基。$\omega$ 处处非零这一事实意味着
\begin{align*}
\Theta_p(E_1,\dots,E_{n-1})=\omega_p(N_p,E_1,\dots,E_{n-1})\ne0.
\end{align*}
由于 $\Theta_p(E_1,\dots,E_{n-1})>0$ 当且仅当 $\omega_p(N_p,E_1,\dots,E_{n-1})>0$，由 $\Theta$ 确定的定向正是命题陈述中所定义的那个定向。
\end{proof}

\begin{example}\label{example:ch15-sphere-orientation}
球面 $S^n$ 是 $\mathbb{R}^{n+1}$ 中的一个超曲面，向量场 $N=x^i\partial/\partial x^i$ 处处不与它相切，故这个向量场在 $S^n$ 上诱导出一个定向。这表明所有球面都可定向。我们把 $S^n$ 的\emph{标准定向}定义为由 $N$ 确定的那个定向。除非另有说明，我们总使用这个定向。（$S^0$ 上的标准定向是给点 $+1\in S^0$ 指派定向 $+1$、给点 $-1\in S^0$ 指派 $-1$ 的那个定向。）
\end{example}

并非每个超曲面都容许一个处处不相切的向量场。（见习题\ref{exercise:ch15-06}。）然而下面的命题给出一个在许多情形下都成立的充分条件。

\begin{proposition}\label{proposition:ch15-regular-level-orientable}
设 $M$ 是一个有向光滑流形，$S\subseteq M$ 是光滑函数 $f:M\to\mathbb{R}$ 的一个正则水平集。则 $S$ 可定向。
\end{proposition}

\begin{proof}
在 $M$ 上任意选取一个 Riemann 度量，令 $N=\operatorname{grad}f|_S$。这些假设蕴含 $N$ 是沿 $S$ 的一个处处不相切的向量场，故结果由命题\ref{proposition:ch15-hypersurface-orientation}得到。
\end{proof}

\subsection{边界的定向}

若 $M$ 是一个带边光滑流形，则 $\partial M$ 是 $M$ 中的一个嵌入超曲面（定理\ref{theorem:ch05-boundary-hypersurface}），而习题\ref{exercise:ch08-04}表明沿 $\partial M$ 总存在一个光滑的向外向量场。由于一个向外向量场处处不与 $\partial M$ 相切，它在 $\partial M$ 上确定一个定向。

\begin{proposition}[边界上的诱导定向]\label{proposition:ch15-boundary-orientation}
设 $M$ 是一个有向光滑带边 $n$ 流形，$n\geqslant1$。则 $\partial M$ 可定向，且沿 $\partial M$ 的所有向外向量场在 $\partial M$ 上确定同一个定向。
\end{proposition}

\begin{remark}
由任意向外向量场所确定的 $\partial M$ 上的定向，称为 $\partial M$ 上的\emph{诱导定向}或\emph{Stokes 定向}。（选用第二个术语，是因为这个定向在第~16 章将讨论的 Stokes 定理中所起的作用。）
\end{remark}

\begin{proof}
令 $n=\dim M$，$\omega$ 是 $M$ 的一个定向形式，$N$ 是沿 $\partial M$ 的一个光滑向外向量场。由命题\ref{proposition:ch15-hypersurface-orientation}，$(n-1)$-形式 $\iota_{\partial M}^*(N\lrcorner\omega)$ 是 $\partial M$ 的一个定向形式，故 $\partial M$ 可定向。

为证这个定向与 $N$ 的选取无关，任取 $p\in\partial M$，在 $p$ 的一个邻域上取 $M$ 的光滑边界坐标 $(x^i)$。若 $N,\widetilde N$ 是沿 $\partial M$ 的两个不同的向外向量场，命题\ref{proposition:ch05-boundary-coords}表明最后一个分量 $N^n(p)$ 与 $\widetilde N^n(p)$ 都为负。$(N_p,\partial/\partial x^1|_p,\dots,\partial/\partial x^{n-1}|_p)$ 与 $(\widetilde N_p,\partial/\partial x^1|_p,\dots,\partial/\partial x^{n-1}|_p)$ 都是 $T_pM$ 的基，它们之间的过渡矩阵行列式等于 $N^n(p)/\widetilde N^n(p)>0$。因此这两个基为 $T_pM$ 确定相同的定向，故 $N$ 与 $\widetilde N$ 为 $T_p\partial M$ 确定相同的定向。（当 $n=1$ 时，所讨论的基就是 $(N_p)$ 与 $(\widetilde N_p)$，它们因都是 $\partial/\partial x^1|_p$ 的负倍数而确定相同定向。）
\end{proof}

\begin{example}\label{example:ch15-sphere-boundary}
这个命题给出了 $S^n$ 可定向的一个更简单的证明，因为它是闭单位球的边界。如此在 $S^n$ 上诱导的定向就是标准定向，你可以自行验证。
\end{example}

\begin{example}\label{example:ch15-halfspace}
让我们确定当 $\mathbb{H}^n$ 本身带有从 $\mathbb{R}^n$ 继承的标准定向时，$\partial\mathbb{H}^n$ 上的诱导定向。我们可在对应
\begin{align*}
(x^1,\dots,x^{n-1},0)\leftrightarrow(x^1,\dots,x^{n-1})
\end{align*}
下把 $\partial\mathbb{H}^n$ 等同于 $\mathbb{R}^{n-1}$。由于向量场 $-\partial/\partial x^n$ 沿 $\partial\mathbb{H}^n$ 向外，$\mathbb{R}^{n-1}$ 的标准坐标标架对 $\partial\mathbb{H}^n$ 正定向当且仅当
\begin{align*}
[-\partial/\partial x^n,\partial/\partial x^1,\dots,\partial/\partial x^{n-1}]
\end{align*}
是 $\mathbb{R}^n$ 的标准定向。这个定向满足
\begin{align*}
[-\partial/\partial x^n,\partial/\partial x^1,\dots,\partial/\partial x^{n-1}]
=-[\partial/\partial x^n,\partial/\partial x^1,\dots,\partial/\partial x^{n-1}]
=(-1)^n[\partial/\partial x^1,\dots,\partial/\partial x^{n-1},\partial/\partial x^n].
\end{align*}
因此，当 $n$ 为偶数时，$\partial\mathbb{H}^n$ 上的诱导定向等于 $\mathbb{R}^{n-1}$ 上的标准定向；而当 $n$ 为奇数时，它与标准定向相反。特别地，$\partial\mathbb{H}^n\cong\mathbb{R}^{n-1}$ 上的标准坐标正定向当且仅当 $n$ 为偶数。（这个事实在第~16 章 Stokes 定理的证明中将很重要。）
\end{example}

\begin{figure}[htbp]
\centering
\includegraphics[width=0.85\textwidth]{fig15-6-boundary.png}
\caption{边界参数化的定向判别法}
\label{figure:ch15-boundary}
\end{figure}

对许多目的而言，描述子流形最有用的方式是借助局部参数化。下一个引理给出一个有用的判别法，用于检查一个边界的局部参数化是否保定向。

\begin{lemma}\label{lemma:ch15-boundary-param}
设 $M$ 是一个有向光滑带边 $n$ 流形。设 $U\subseteq\mathbb{R}^{n-1}$ 是开集，$a,b$ 是满足 $a<b$ 的实数，$F:(a,b]\times U\to M$ 是一个光滑嵌入，它限制为从 $\{b\}\times U$ 到 $\partial M$ 的一个嵌入。则由 $f(x)=F(b,x)$ 给出的参数化 $f:U\to\partial M$ 对 $\partial M$ 保定向，当且仅当 $F$ 对 $M$ 保定向。
\end{lemma}

\begin{proof}
令 $x$ 是 $U$ 的任意一点，令 $p=f(x)=F(b,x)\in\partial M$（图\ref{figure:ch15-boundary}）。$F$ 是嵌入这一假设意味着线性映射 $\mathrm{d}F_{(b,x)}:(T_b\mathbb{R}\oplus T_x\mathbb{R}^{n-1})\to T_pM$ 是双射。由于 $\mathrm{d}F_{(b,x)}$ 在 $T_x\mathbb{R}^{n-1}$ 上的限制等于 $\mathrm{d}f_x:T_x\mathbb{R}^{n-1}\to T_p\partial M$（它已经是单射），可知
\begin{align*}
\mathrm{d}F\left(\partial/\partial s|_{(b,x)}\right)\notin T_p\partial M,
\end{align*}
其中 $s$ 表示 $(a,b]$ 上的坐标。

由 $\gamma(t)=F(b-t,x)$ 定义一条光滑曲线 $\gamma:[0,\varepsilon)\to M$。这条曲线满足 $\gamma(0)=p$，且 $\gamma'(0)=-\mathrm{d}F(\partial/\partial s|_{(b,x)})\notin T_p\partial M$。由此可知 $-\mathrm{d}F(\partial/\partial s|_{(b,x)})$ 是向内指向的，因而 $\mathrm{d}F(\partial/\partial s|_{(b,x)})$ 是向外指向的。

诱导定向的定义给出如下等价关系：
\begin{align*}
F\text{ 对 }M\text{ 保定向}
&\iff\bigl(\mathrm{d}F(\partial/\partial s),\mathrm{d}F(\partial/\partial x^1),\dots,\mathrm{d}F(\partial/\partial x^{n-1})\bigr)\text{ 对 }TM\text{ 定向}\\
&\iff\bigl(\mathrm{d}F(\partial/\partial x^1),\dots,\mathrm{d}F(\partial/\partial x^{n-1})\bigr)\text{ 对 }T\partial M\text{ 定向}\\
&\iff\bigl(\mathrm{d}f(\partial/\partial x^1),\dots,\mathrm{d}f(\partial/\partial x^{n-1})\bigr)\text{ 对 }T\partial M\text{ 定向}\\
&\iff f\text{ 对 }\partial M\text{ 保定向}.
\end{align*}
\end{proof}

\begin{figure}[htbp]
\centering
\includegraphics[width=0.8\textwidth]{fig15-7-spherecoord.png}
\caption{借助球坐标参数化球面}
\label{figure:ch15-spherecoord}
\end{figure}

下面说明这个引理如何使用。

\begin{example}\label{example:ch15-spherical-param}
球坐标（例\ref{example:app-spherical}）如下给出 $S^2$ 的一个光滑局部参数化。令 $U$ 为开矩形 $(0,\pi)\times(0,2\pi)\subseteq\mathbb{R}^2$，令 $X:U\to\mathbb{R}^3$ 为映射
\begin{align*}
X(\varphi,\theta)=(\sin\varphi\cos\theta,\sin\varphi\sin\theta,\cos\varphi)
\end{align*}
（图\ref{figure:ch15-spherecoord}）。我们可以利用它是由 $F:(0,1]\times U\to\overline{\mathbb{B}}^3$（定义为
\begin{align*}
F(\rho,\varphi,\theta)=(\rho\sin\varphi\cos\theta,\rho\sin\varphi\sin\theta,\rho\cos\varphi)
\end{align*}
给出的三维球坐标参数化）的限制这一事实，来检查 $X$ 保定向还是反定向。由于 $F(1,\varphi,\theta)=X(\varphi,\theta)$，引理\ref{lemma:ch15-boundary-param}的假设满足。直接计算，$F$ 的 Jacobian 行列式为 $\rho^2\sin\varphi$，它在 $(0,1]\times U$ 上为正。由引理\ref{lemma:ch15-boundary-param}，$X$ 保定向。
\end{example}
\end{document}
